% Options for packages loaded elsewhere
\PassOptionsToPackage{unicode}{hyperref}
\PassOptionsToPackage{hyphens}{url}
%
\documentclass[
]{article}
\usepackage{amsmath,amssymb}
\usepackage{lmodern}
\usepackage{iftex}
\ifPDFTeX
  \usepackage[T1]{fontenc}
  \usepackage[utf8]{inputenc}
  \usepackage{textcomp} % provide euro and other symbols
\else % if luatex or xetex
  \usepackage{unicode-math}
  \defaultfontfeatures{Scale=MatchLowercase}
  \defaultfontfeatures[\rmfamily]{Ligatures=TeX,Scale=1}
\fi
% Use upquote if available, for straight quotes in verbatim environments
\IfFileExists{upquote.sty}{\usepackage{upquote}}{}
\IfFileExists{microtype.sty}{% use microtype if available
  \usepackage[]{microtype}
  \UseMicrotypeSet[protrusion]{basicmath} % disable protrusion for tt fonts
}{}
\makeatletter
\@ifundefined{KOMAClassName}{% if non-KOMA class
  \IfFileExists{parskip.sty}{%
    \usepackage{parskip}
  }{% else
    \setlength{\parindent}{0pt}
    \setlength{\parskip}{6pt plus 2pt minus 1pt}}
}{% if KOMA class
  \KOMAoptions{parskip=half}}
\makeatother
\usepackage{xcolor}
\usepackage{longtable,booktabs,array}
\usepackage{calc} % for calculating minipage widths
% Correct order of tables after \paragraph or \subparagraph
\usepackage{etoolbox}
\makeatletter
\patchcmd\longtable{\par}{\if@noskipsec\mbox{}\fi\par}{}{}
\makeatother
% Allow footnotes in longtable head/foot
\IfFileExists{footnotehyper.sty}{\usepackage{footnotehyper}}{\usepackage{footnote}}
\makesavenoteenv{longtable}
\setlength{\emergencystretch}{3em} % prevent overfull lines
\providecommand{\tightlist}{%
  \setlength{\itemsep}{0pt}\setlength{\parskip}{0pt}}
\setcounter{secnumdepth}{-\maxdimen} % remove section numbering
\ifLuaTeX
  \usepackage{selnolig}  % disable illegal ligatures
\fi
\IfFileExists{bookmark.sty}{\usepackage{bookmark}}{\usepackage{hyperref}}
\IfFileExists{xurl.sty}{\usepackage{xurl}}{} % add URL line breaks if available
\urlstyle{same} % disable monospaced font for URLs
\hypersetup{
  hidelinks,
  pdfcreator={LaTeX via pandoc}}

\author{}
\date{}

\begin{document}

\begin{longtable}[]{@{}
  >{\raggedright\arraybackslash}p{(\columnwidth - 0\tabcolsep) * \real{1.0000}}@{}}
\toprule()
\begin{minipage}[b]{\linewidth}\raggedright
UCS503P --- Project Proposal

\begin{quote}
\textbf{Campus Sports Equipment Issue--Return System}

\emph{A Digital-Credential Platform Replacing Physical ID-Card
Collateral for Campus Sports Equipment Lending}\\
Project Proposal for UCS503P
\end{quote}

\begin{longtable}[]{@{}
  >{\raggedright\arraybackslash}p{(\columnwidth - 2\tabcolsep) * \real{0.5000}}
  >{\raggedright\arraybackslash}p{(\columnwidth - 2\tabcolsep) * \real{0.5000}}@{}}
\toprule()
\begin{minipage}[b]{\linewidth}\raggedright
\begin{quote}
\textbf{Student Name}
\end{quote}
\end{minipage} & \begin{minipage}[b]{\linewidth}\raggedright
\begin{quote}
Keshav, Manan ,Tushar
\end{quote}
\end{minipage} \\
\midrule()
\endhead
\begin{minipage}[t]{\linewidth}\raggedright
\begin{quote}
\textbf{Roll Number}
\end{quote}
\end{minipage} & \begin{minipage}[t]{\linewidth}\raggedright
\begin{quote}
1024240092 ,1024240136 ,1024240153
\end{quote}
\end{minipage} \\
\begin{minipage}[t]{\linewidth}\raggedright
\begin{quote}
\textbf{Branch}
\end{quote}
\end{minipage} & \begin{minipage}[t]{\linewidth}\raggedright
\begin{quote}
Artificial Intelligence and Machine Learning (AIML)
\end{quote}
\end{minipage} \\
\begin{minipage}[t]{\linewidth}\raggedright
\begin{quote}
\textbf{Institute}
\end{quote}
\end{minipage} & \begin{minipage}[t]{\linewidth}\raggedright
\begin{quote}
Thapar Institute of Engineering and Technology (TIET), Patiala
\end{quote}
\end{minipage} \\
\begin{minipage}[t]{\linewidth}\raggedright
\begin{quote}
\textbf{Course}
\end{quote}
\end{minipage} & \begin{minipage}[t]{\linewidth}\raggedright
\begin{quote}
UCS503
\end{quote}
\end{minipage} \\
\bottomrule()
\end{longtable}

\begin{quote}
\textbf{1Higher-order goal}

This project aims to improve student welfare and the operational
efficiency of campus sports infrastructure by removing the
administrative friction built into the current equipment-lending
process. While the immediate trigger is a single policy at one college
--- students must deposit their college ID card to borrow sports
equipment --- the underlying engineering objective is broader: design a
deployable, reusable digital-credential and inventory-tracking system
that can generalize to any campus asset-lending scenario (sports
equipment today; library books or lab equipment tomorrow).

\textbf{2Time-to-value}

•We will deliver meaningful increments quickly by prototyping the core
issue--return workflow (digital token generation and staff-side
scanning) and validating it with real students and sports-office staff
within a short iteration window.

•We will prioritize fast feedback over perfection by using weekly
iterations and CI-backed automation early, to keep release friction low.

•Success is defined by what can be delivered and measured in the first
iteration --- a working issue/return loop with timestamps --- then
improved based on user feedback.

\textbf{3Problem Statement}

At present, students who want to borrow sports equipment (bats, rackets,
footballs, etc.) from the college sports office must hand over their
physical college ID card as collateral, which is returned only when the
equipment is returned. In practice this creates real, recurring
friction:

•Access barrier: students who forget their ID card cannot borrow
equipment at all, and often end up borrowing a friend\textquotesingle s
card instead --- which breaks accountability and misattributes usage.

•No visibility: there is no way to check equipment availability before
walking to the sports office, so students frequently make a wasted trip
and return empty-handed.
\end{quote}

1\strut
\end{minipage} \\
\midrule()
\endhead
\bottomrule()
\end{longtable}

\begin{longtable}[]{@{}
  >{\raggedright\arraybackslash}p{(\columnwidth - 0\tabcolsep) * \real{1.0000}}@{}}
\toprule()
\begin{minipage}[b]{\linewidth}\raggedright
UCS503P --- Project Proposal

\begin{quote}
•Manual, paper-based tracking: issue and return records are kept in
physical registers, making it hard to trace who has what, for how long,
or in what condition.

•Weak enforcement: the only real penalty mechanism today is withholding
the ID card itself, which is a blunt, all-or-nothing control rather than
a graded, trackable consequence such as a fine.
\end{quote}

\begin{longtable}[]{@{}
  >{\raggedright\arraybackslash}p{(\columnwidth - 2\tabcolsep) * \real{0.5000}}
  >{\raggedright\arraybackslash}p{(\columnwidth - 2\tabcolsep) * \real{0.5000}}@{}}
\toprule()
\begin{minipage}[b]{\linewidth}\raggedright
\textbf{Current Process}
\end{minipage} & \begin{minipage}[b]{\linewidth}\raggedright
\textbf{Proposed System}
\end{minipage} \\
\midrule()
\endhead
\begin{minipage}[t]{\linewidth}\raggedright
\begin{quote}
ID card held as physical collateral
\end{quote}
\end{minipage} & Time-bound digital token; card stays with student \\
\begin{minipage}[t]{\linewidth}\raggedright
\begin{quote}
No visibility into availability before visiting
\end{quote}
\end{minipage} & \begin{minipage}[t]{\linewidth}\raggedright
\begin{quote}
Live availability + slot reservation
\end{quote}
\end{minipage} \\
\begin{minipage}[t]{\linewidth}\raggedright
\begin{quote}
Paper register, hard to trace history
\end{quote}
\end{minipage} & \begin{minipage}[t]{\linewidth}\raggedright
\begin{quote}
Digital issue--return log with timestamps
\end{quote}
\end{minipage} \\
\begin{minipage}[t]{\linewidth}\raggedright
\begin{quote}
Card withheld as sole penalty
\end{quote}
\end{minipage} & \begin{minipage}[t]{\linewidth}\raggedright
\begin{quote}
Graded fines tied to no-dues clearance
\end{quote}
\end{minipage} \\
\bottomrule()
\end{longtable}

\begin{quote}
\textbf{Why this matters now (contextual relevance):} A large college
with many sports offerings has correspondingly heavy equipment turnover;
even a modest reduction in wasted trips and untracked loans meaningfully
improves both student experience and sports-office accountability.

\textbf{4Proposed Solution}

\textbf{4.1Overview}

Build a mobile-friendly web platform (Progressive Web App) that replaces
the physical ID-card deposit with a time-bound, non-transferable digital
credential, while borrowing the accountability model of a library
issue--return system. Core components:

•Student portal: authenticate with college ID, browse live equipment
availability by category and location, optionally reserve a time slot.

•Digital issue credential: a QR/token generated per request, bound to
the student, the specific item, and a time slot --- not a reusable
general-purpose ID.

•Staff scanning app: a lightweight interface for sports-office staff to
scan the token at issue and again at return, replacing the paper
register.

•Fine engine: automatic fine calculation for late or damaged returns,
feeding into the college\textquotesingle s existing "no dues" clearance
process at semester/degree end.

•Admin dashboard: inventory status, backlog, utilization, and
most-borrowed items for the sports office.

\textbf{4.2Core workflow}

•Student logs in with their college ID and browses live equipment
availability, optionally reserving a slot.

•System generates a time-bound digital token for that student, item, and
slot --- the physical ID card is no longer required.

•Staff scan the token at the counter; the item is marked Issued, the due
time starts, and the student keeps their ID card.

•On return, staff scan the same request again (or a return-side token)
and record the item\textquotesingle s condition; the item is marked
Returned or Damaged.
\end{quote}

2
\end{minipage} \\
\midrule()
\endhead
\bottomrule()
\end{longtable}

\begin{longtable}[]{@{}
  >{\raggedright\arraybackslash}p{(\columnwidth - 0\tabcolsep) * \real{1.0000}}@{}}
\toprule()
\begin{minipage}[b]{\linewidth}\raggedright
UCS503P --- Project Proposal

\begin{quote}
•If the return is late or the item is damaged, a fine is automatically
logged against the student\textquotesingle s record and surfaced on both
the student dashboard and the no-dues clearance system.

\textbf{4.3Operational constraints for an educational, tier-2/3 setting}

•Keep the solution usable on low-to-mid range devices via a responsive
Progressive Web App rather than a native app, avoiding install friction.

•Avoid dependence on advanced research algorithms; focus on workflow
correctness, credential design, and accountability.

•Design for deployability using common web hosting and managed
databases, with a hybrid (app + manual fallback) rollout during the
pilot so the sports office is never fully blocked by a system outage.

\textbf{4.4Extension: Facility \& Court Slot Booking}\\
•The complex also has bookable facilities ---
badminton/tennis/basketball courts, the gym, the pool --- not just loose
equipment. Since the reservation engine in Section 4.2 already
generalizes to "a student, an asset, a time slot," the same platform is
extended to let students book a facility time slot directly, instead of
only equipment:

•A facility (e.g. "Badminton Court 2") exposes bookable hourly slots,
shown with live availability alongside equipment, using the same
browse-and-reserve screen.

•Double-booking is prevented with a simple database-level uniqueness\\
constraint/transaction on (facility, slot) at confirmation time ---
deliberately simpler than a distributed lock (e.g. Redis), which is
unnecessary infrastructure at this scale and would work against the
course\textquotesingle s "avoid advanced/research-heavy infrastructure"
guidance.

•No online payment or membership tier is introduced: access stays
governed by the student\textquotesingle s existing college ID rather
than a paid subscription, keeping this an extension of the same
credential system rather than a separate commercial booking product.

•This is scoped as a later-iteration, opt-in module (see Section 9.2)
built on the existing token/reservation engine, not a parallel system
--- it is what makes the platform a genuine "sports complex" system
rather than an equipment-only tool.

\textbf{5Solution Approach}

This is an engineering course project, so the focus is on credential
design, workflow correctness, and system reliability rather than novel
algorithm research.

\textbf{5.1Digital credential design}

•Each token is time-bound, non-transferable, and scoped to one student,
one item, and one slot --- unlike a general-purpose reusable ID.
\end{quote}

3\strut
\end{minipage} \\
\midrule()
\endhead
\bottomrule()
\end{longtable}

\begin{longtable}[]{@{}
  >{\raggedright\arraybackslash}p{(\columnwidth - 0\tabcolsep) * \real{1.0000}}@{}}
\toprule()
\begin{minipage}[b]{\linewidth}\raggedright
UCS503P --- Project Proposal

\begin{quote}
•Unlike a single-scan canteen-style token, the credential is two-state
(Issue scan, then Return scan), since equipment lending has two distinct
events rather than one completed transaction.

•A manual staff-side override (ID lookup) is retained as a fallback for
scan failures or network issues, so the desk is never blocked.

\textbf{5.2Web architecture}

A typical 3-tier web architecture:

•Frontend: responsive PWA for students (browse, reserve, view token) and
a simplified scanning view for staff.

•Backend API: authentication, inventory queries, token
issuance/validation, status transitions, and fine calculation.

•Database: students, equipment items, batches/units, issue-return
history, fines, and notifications.

\textbf{5.3CI/CD and fast delivery enabler}

•Automatically build and run tests on every change.

•Produce repeatable deployment artifacts to staging.

•Enable safe and frequent releases of incremental features.

\textbf{6Evaluation Criterion \ldots{} measurable and attributable}

We define success using measurable metrics tied to the core workflow.

\textbf{6.1Primary evaluation metric}

\textbf{Time-to-Issue (TTI):} the time between a student generating a
digital token and receiving the physical equipment at the counter.

•Target: median TTI ≤ 5 minutes during the pilot window.

•Attribution: measured from database timestamps (token generation vs.
staff issue-scan).

\textbf{6.2Secondary evaluation metrics}

•Empty-handed rate: percentage of visits where a student could not
obtain equipment due to unavailability, compared against a baseline
collected before the pilot.

•Return-on-time rate: percentage of items returned before their due
time.

•Fine-to-resolution rate: percentage of issued fines cleared before
semester-end dues clearance.

•Staff processing time: time taken per issue/return transaction compared
with the existing paper-register process.
\end{quote}

•User clarity: student-reported clarity/satisfaction using a simple 1--5
feedback question.

4
\end{minipage} \\
\midrule()
\endhead
\bottomrule()
\end{longtable}

\begin{longtable}[]{@{}
  >{\raggedright\arraybackslash}p{(\columnwidth - 0\tabcolsep) * \real{1.0000}}@{}}
\toprule()
\begin{minipage}[b]{\linewidth}\raggedright
UCS503P --- Project Proposal

\begin{quote}
\textbf{6.3Pilot validation plan}

•Recruit 10--20 pilot users (students) and 2--3 sports-office staff for
a two-week pilot run alongside the existing manual process.

•Compare metrics before/after within the iteration window by logging
baseline estimates via short interviews and existing register data.

\textbf{7Scalability}

The proposition should scale in both theoretical and practical senses.

\textbf{Scalable (theoretically):} the system should support growth in
the number of equipment items, concurrent student requests, and
dashboard usage by:

•decoupling components (frontend/backend),

•enabling pagination and indexing for common queries (availability
lookups, student history),

•using stateless API design so multiple sports-office counters can
operate concurrently.

\textbf{Operationally deployable (practically):} the system should run
on common web hosting:

•managed database,

•containerized or platform-hosted deployment,

•CI/CD pipeline for staging/production.

Because the credential-and-inventory engine is generic (student, asset,
slot, condition, fine), it is designed to extend beyond sports equipment
to other campus asset-lending contexts --- library books, lab equipment,
or event gear --- without a redesign, which is the
project\textquotesingle s primary scalability argument.

\textbf{8Engine availability heuristic}

A mature set of "engines" (tooling/services) is available to implement
this as a web-based project; indicative choices are listed below (final
stack to be confirmed at design time):

•Frontend: React (or Next.js) as an installable Progressive Web App, for
the student portal and staff scanning view.

•Backend API: Node.js/Express or FastAPI, layered as Route → Service →
Repository.

•Database: PostgreSQL (managed instance) via an ORM (Prisma/SQLAlchemy)
for schema migrations.

•Auth: JWT-based sessions, with an integration point for college SSO/ID
verification where available.

•QR code generation (e.g. qrcode) and browser-based camera scanning
(e.g. html5-qrcode) library.
\end{quote}

5
\end{minipage} \\
\midrule()
\endhead
\bottomrule()
\end{longtable}

\begin{longtable}[]{@{}
  >{\raggedright\arraybackslash}p{(\columnwidth - 0\tabcolsep) * \real{1.0000}}@{}}
\toprule()
\begin{minipage}[b]{\linewidth}\raggedright
UCS503P --- Project Proposal

\begin{quote}
•Notification channel for due-time reminders (in-app first via web push;
email/SMS optional later).

•Test framework: Jest/PyTest for unit and integration tests.

•CI/CD: GitHub Actions runner, deploying containerized builds to a
staging environment. We will use these mature components to focus on
workflow design, credential correctness, and measurement rather than
inventing new infrastructure.

\textbf{9Project scope and deliverables}

\textbf{9.1Initial deliverable}\\
•Student login and live equipment browsing\\
•Digital token generation for a request\\
•Staff scanning app: issue-side scan and status update •Basic logging
and timestamps for TTI measurement •CI: build + backend unit tests +
linting\\
•CD to staging with a single command or pipeline job

\textbf{9.2Subsequent deliverables}\\
•Return-side scan with condition capture (Returned / Damaged) and
automatic fine calculation\\
•Live slot booking/reservation to prevent wasted trips\\
•Due-time reminder notifications (in-app, then email/SMS)\\
•No-dues integration or export for semester-end clearance\\
•Admin dashboard for utilization, backlog, and most-borrowed items\\
•Improved search/filtering and indexed queries for scale readiness

\textbf{9.3Indicative iteration timeline}
\end{quote}

\begin{longtable}[]{@{}
  >{\raggedright\arraybackslash}p{(\columnwidth - 2\tabcolsep) * \real{0.5000}}
  >{\raggedright\arraybackslash}p{(\columnwidth - 2\tabcolsep) * \real{0.5000}}@{}}
\toprule()
\begin{minipage}[b]{\linewidth}\raggedright
\textbf{Timeframe}
\end{minipage} & \begin{minipage}[b]{\linewidth}\raggedright
\textbf{Focus}
\end{minipage} \\
\midrule()
\endhead
\begin{minipage}[t]{\linewidth}\raggedright
\begin{quote}
Week 1--2
\end{quote}
\end{minipage} & \begin{minipage}[t]{\linewidth}\raggedright
\begin{quote}
Setup, data model, student login, CI pipeline scaffolding
\end{quote}
\end{minipage} \\
\begin{minipage}[t]{\linewidth}\raggedright
\begin{quote}
Week 3--4
\end{quote}
\end{minipage} & \begin{minipage}[t]{\linewidth}\raggedright
\begin{quote}
Live availability browsing, digital token generation, staff issue-scan
\end{quote}
\end{minipage} \\
\begin{minipage}[t]{\linewidth}\raggedright
\begin{quote}
Week 5--6
\end{quote}
\end{minipage} & \begin{minipage}[t]{\linewidth}\raggedright
\begin{quote}
Return-scan, condition capture, automatic fine engine
\end{quote}
\end{minipage} \\
\begin{minipage}[t]{\linewidth}\raggedright
\begin{quote}
Week 7--8
\end{quote}
\end{minipage} & \begin{minipage}[t]{\linewidth}\raggedright
\begin{quote}
Slot booking/reservation, due-time notifications
\end{quote}
\end{minipage} \\
\begin{minipage}[t]{\linewidth}\raggedright
\begin{quote}
Week 9--10
\end{quote}
\end{minipage} & \begin{minipage}[t]{\linewidth}\raggedright
\begin{quote}
Admin dashboard, no-dues export, pilot with real users
\end{quote}
\end{minipage} \\
\begin{minipage}[t]{\linewidth}\raggedright
\begin{quote}
Week 11--12
\end{quote}
\end{minipage} & \begin{minipage}[t]{\linewidth}\raggedright
\begin{quote}
Feedback-driven fixes, scale-readiness pass, final evaluation
\end{quote}
\end{minipage} \\
\bottomrule()
\end{longtable}

\begin{quote}
\textbf{10Risks and mitigations}

•Authentication dependency: the system assumes access to a college
SSO/ID-verification mechanism; if unavailable, fall back to admin-issued
student credentials for the pilot.
\end{quote}

6\strut
\end{minipage} \\
\midrule()
\endhead
\bottomrule()
\end{longtable}

\begin{longtable}[]{@{}
  >{\raggedright\arraybackslash}p{(\columnwidth - 0\tabcolsep) * \real{1.0000}}@{}}
\toprule()
\begin{minipage}[b]{\linewidth}\raggedright
UCS503P --- Project Proposal

\begin{quote}
•Staff adoption: paper registers are familiar and trusted; mitigate with
a simple scanning UI, a hybrid rollout alongside the manual process, and
a manual-entry fallback for scan failures.

•Damage/loss disputes: mandate a condition note (and optionally a photo)
at both issue and return to avoid disagreements about pre-existing
damage.

•No-dues integration access: the college\textquotesingle s
degree-clearance system may not expose an API; mitigate by exporting a
fine list the sports office can reconcile manually during the pilot,
with deeper integration explored later.

•Student adoption resistance: run a short awareness/onboarding push at
pilot launch, since the workflow changes a habit students already know.

\textbf{11Summary}

This project proposes a deployable web platform that removes the
ID-card-as-collateral requirement from campus sports equipment lending,
replacing it with a time-bound digital credential modelled on library
issue--return accountability. It meets the course criteria by
emphasizing time-to-value through rapid iterations, implementing CI/CD
to support frequent safe releases, and using measurable evaluation
criteria --- especially time-to-issue and empty-handed-visit reduction.
It remains focused on engineering workflow, credential design, and
scalability readiness (toward other campus asset-lending use cases)
rather than research-driven novelty.
\end{quote}

7
\end{minipage} \\
\midrule()
\endhead
\bottomrule()
\end{longtable}

\end{document}
